\section{Algebraic structures}
\label{section-algebraic-structures}

\noindent
In this section we mildly formalize the notions we have
encountered in the sections above.

\begin{definition}
\label{definition-algebraic-structure}
A {\it type of algebraic structure} is given by a category $\mathcal{C}$
and a functor $F : \mathcal{C} \to \textit{Sets}$ with the
following properties
\begin{enumerate}
\item $F$ is faithful,
\item $\mathcal{C}$ has limits and $F$ commutes with limits,
\item $\mathcal{C}$ has filtered colimits and $F$ commutes with them, and
\item $F$ reflects isomorphisms.
\end{enumerate}
\end{definition}

\noindent
We make this definition to point out the properties
we will use in a number of arguments below.  But we will
not actually study this notion in any great detail, since we
are prohibited from studying ``big'' categories by convention,
except for those listed in
Categories, Remark \ref{categories-remark-big-categories}.
Among those the following have the required properties.

\begin{lemma}
\label{lemma-list-algebraic-structures}
The following categories, endowed with the obvious forgetful
functor, define types of algebraic structures:
\begin{enumerate}
\item The category of pointed sets.
\item The category of abelian groups.
\item The category of groups.
\item The category of monoids.
\item The category of rings.
\item The category of $R$-modules for a fixed ring $R$.
\item The category of Lie algebras over a fixed field.
\end{enumerate}
\end{lemma}

\begin{proof}
Omitted.
\end{proof}

\noindent
From now on we will think of a (pre)sheaf of algebraic structures
and their stalks, in terms of the underlying (pre)sheaf of sets.
This is allowable by Lemmas \ref{lemma-sheaves-structure} and
\ref{lemma-stalk-presheaf-values-in-category}.

\medskip\noindent
In the rest of this section we point out some results
on algebraic structures that will be useful in the future.

\begin{lemma}
\label{lemma-properties-algebraic-structures}
Let $(\mathcal{C}, F)$ be a type of algebraic structure.
\begin{enumerate}
\item $\mathcal{C}$ has a final object $0$ and $F(0) = \{ * \}$.
\item $\mathcal{C}$ has products and $F(\prod A_i) = \prod F(A_i)$.
\item $\mathcal{C}$ has fibre products and
$F(A \times_B C) = F(A)\times_{F(B)}F(C)$.
\item $\mathcal{C}$ has equalizers, and if $E \to A$
is the equalizer of $a, b : A \to B$, then
$F(E) \to F(A)$ is the equalizer of $F(a), F(b) : F(A) \to F(B)$.
\item $A \to B$ is a monomorphism if and only if
$F(A) \to F(B)$ is injective.
\item if $F(a) : F(A) \to F(B)$ is surjective, then
$a$ is an epimorphism.
\item given $A_1 \to A_2 \to A_3 \to \ldots$, then
$\colim A_i$ exists and $F(\colim A_i) = \colim F(A_i)$,
and more generally for any filtered colimit.
\end{enumerate}
\end{lemma}

\begin{proof}
Omitted. The only interesting statement is (5) which
follows because $A \to B$ is a monomorphism if and only if
$A \to A \times_B A$ is an isomorphism, and then applying
the fact that $F$ reflects isomorphisms.
\end{proof}

\begin{lemma}
\label{lemma-image-contained-in}
Let $(\mathcal{C}, F)$ be a type of algebraic structure.
Suppose that $A, B, C \in \Ob(\mathcal{C})$.
Let $f : A \to B$ and $g : C \to B$ be morphisms of
$\mathcal{C}$. If $F(g)$ is injective, and
$\Im(F(f)) \subset \Im(F(g))$, then
$f$ factors as $f = g \circ t$ for some morphism
$t : A \to C$.
\end{lemma}

\begin{proof}
Consider $A \times_B C$. The assumptions imply that
$F(A \times_B C) = F(A) \times_{F(B)} F(C) = F(A)$.
Hence $A = A \times_B C$ because $F$ reflects isomorphisms.
The result follows.
\end{proof}

\begin{example}
\label{example-application-lemma-image-contained-in}
The lemma will be applied often to the following situation.
Suppose that we have a diagram
$$
\xymatrix{
A \ar[r] & B \ar[d] \\
C \ar[r] & D
}
$$
in $\mathcal{C}$. Suppose $C \to D$ is injective on underlying
sets, and suppose that the composition $A \to B \to D$ has
image on underlying sets in the image of $C \to D$.
Then we get a commutative diagram
$$
\xymatrix{
A \ar[r] \ar[d] & B \ar[d] \\
C \ar[r] & D
}
$$
in $\mathcal{C}$.
\end{example}

\begin{example}
\label{example-sheaf-product-pointwise-algebraic-structure}
Let $F : \mathcal{C} \to \textit{Sets}$ be a
type of algebraic structures.
Let $X$ be a topological space. Suppose that
for every $x \in X$ we are given an object
$A_x \in \Ob(\mathcal{C})$. Consider the presheaf
$\Pi$ with values in $\mathcal{C}$ on $X$ defined by the rule
$\Pi(U) = \prod_{x \in U} A_x$ (with obvious restriction
mappings). Note that the associated presheaf of sets
$U \mapsto F(\Pi(U)) = \prod_{x \in U} F(A_x)$ is a sheaf
by Example \ref{example-sheaf-product-pointwise}.
Hence $\Pi$ is a sheaf of algebraic structures
of type $(\mathcal{C} , F)$. This gives many examples
of sheaves of abelian groups, groups, rings, etc.
\end{example}
